% ===========================================================
%  THEOREM 1 -- PROOF SKETCH (with gap markers)
%  v10 AAAI 2027
% ===========================================================

\documentclass[11pt,a4paper]{article}
\usepackage[margin=1in]{geometry}
\usepackage{amsmath,amssymb,amsthm,bm,hyperref,booktabs}
\newtheorem{theorem}{Theorem}
\newtheorem{corollary}[theorem]{Corollary}
\newtheorem{lemma}[theorem]{Lemma}
\newtheorem{proposition}[theorem]{Proposition}
\newcommand{\R}{\mathbb{R}}
\newcommand{\E}{\mathbb{E}}
\newcommand{\TODO}[1]{\textbf{[TODO: user verify -- #1]}}
\title{Proof of Theorem~1 (sketch)\\
\large Direction-invariant label flip under linear batch correction}
\date{v10 AAAI 2027 -- 2026-04-24}
\begin{document}
\maketitle

\noindent\textbf{Reading guide.} This document is a \emph{proof sketch}
with explicit gap markers of the form \TODO{...}. Each marker flags a
step that the human author will need to expand with a rigorous
finite-sample or asymptotic argument before the AAAI camera-ready.

\section{High-level strategy}
\begin{enumerate}
\item[(a)] Decompose $X$ into orthogonal label/batch/residual components.
\item[(b)] Apply the linear correction operator $T_B$ in
  \eqref{eq:TB-recap} and quantify how much of the
  label signal is lost together with the batch signal.
\item[(c)] Identify the exact fraction of label signal surviving; show
  the sign of the surviving component flips at $\rho^2=\tfrac12$.
\item[(d)] Carry this through to the empirical risk minimiser
  $w_{\text{trained}}$ via a ridge-consistency argument.
\item[(e)] Translate the sign of $\langle w_{\text{trained}}, w_Y\rangle$
  into the AUC statement.
\end{enumerate}

\section{Step (a): Orthogonal decomposition}

Under Assumption~\ref{asm:additive} (see setup document),
\begin{equation}
  X \;=\; Y\,w_Y^{\!\top}\lVert\mu_1-\mu_0\rVert
     \;+\; \tilde B\,M_B\,w_B^{\!\top}
     \;+\; \mathcal{E},
  \label{eq:decomp}
\end{equation}
where $\tilde B\in\R^{n\times (K-1)}$ is the centred batch indicator and
$M_B\in\R^{(K-1)\times 1}$ aggregates the batch centroids onto the
dominant batch direction. The residual $\mathcal{E}$ has rows that are
independent, mean-zero, and sub-Gaussian by Assumption~\ref{asm:additive}.

\paragraph{Correlation structure.} By construction,
\begin{equation}
  \bigl\langle w_Y,\, w_B \bigr\rangle \;=\; \rho,
  \qquad
  \lVert w_Y\rVert = \lVert w_B\rVert = 1.
  \label{eq:wywb}
\end{equation}
The two directions are not orthogonal precisely because batch membership
predicts the label (Assumption~\ref{asm:conf}).

\section{Step (b): Action of $T_B$}

The linear correction~(\ref{eq:TB-matrix} in the setup) has the compact
form
\begin{equation}
  T_B(X) \;=\; (I-M_B)\,X_{\!B} \;+\; X_{\!Y}^{\text{kept}} \;+\; \mathcal{E}',
  \label{eq:TB-recap}
\end{equation}
where $X_B$ is the batch component of~\eqref{eq:decomp} and
$X_Y^{\text{kept}}$ is the fraction of the label component that the
covariate-preserving construction attempts to retain.

\paragraph{Retention coefficient.} Following the standard ComBat
derivation (Johnson--Li--Rabinovic 2007, equations~2--4) when the
label covariate is included in the OLS design, a fraction of the label
signal is unavoidably projected away. We claim:
\begin{lemma}[Label retention under covariate-preserving ComBat]
\label{lem:retention}
Under Assumptions~\ref{asm:additive}--\ref{asm:n},
\begin{equation}
   T_B(X)\, w_Y \;=\; (1-\rho^{2})\,Y\,\lVert\mu_1-\mu_0\rVert
   \;-\;\rho^{2}\,Y\,\lVert\mu_1-\mu_0\rVert \;+\; O_{\!p}(n^{-1/2}).
   \label{eq:retention}
\end{equation}
Equivalently, projecting $T_B(X)$ onto $w_Y$ leaves a signed coefficient
\(
  c(\rho) \;=\; 1-2\rho^{2}
\)
times the original label amplitude.
\end{lemma}

\paragraph{Proof of Lemma~\ref{lem:retention} (K=2, balanced).}
We carry out the calculation rigorously for two balanced cohorts,
$|\pi_1|=|\pi_2|=n/2$, with the additional assumption that $Y$ has
cross-cohort divergence parameter $\rho_{YY}\in[-1,1]$ measuring the
cosine agreement between within-cohort label-direction estimates. The
Gaussian model is
\begin{equation}
   X_i \;=\; \alpha + \beta_{B_i}^{\!\top} Y_i + \gamma_{B_i} + \varepsilon_i,
   \qquad \beta_{B=k} = \sigma_Y \, e_{Y,k}
   \label{eq:model-plat}
\end{equation}
with $\langle e_{Y,1}, e_{Y,2}\rangle = \rho_{YY}$ and batch-mean
drift $\gamma_k = \sigma_B\,(2k{-}3)\,e_B$, $\langle e_{Y,1},e_B\rangle=\rho$.
Let $Y\in\{0,1\}$ with $\E Y=1/2$ globally but $\E[Y\mid B=1]=q$,
$\E[Y\mid B=2]=1-q$ (confounding parameter).

\paragraph{OLS covariate estimate.}
With design $D=[\mathbf 1_n\;Y\;\mathbf 1_{B=2}]$, the OLS estimate
of the $Y$ coefficient $\hat\beta$ is, to leading order in $1/n$,
\begin{equation}
  \hat\beta \;=\; \frac{\sum_i (Y_i-\bar Y)X_i}{\sum_i(Y_i-\bar Y)^2}.
  \label{eq:betahat}
\end{equation}
Because the true slope differs between cohorts
($\beta_1=\sigma_Y e_{Y,1}$ vs $\beta_2=\sigma_Y e_{Y,2}$), the common
$\hat\beta$ converges to the cohort-size-weighted average
$\hat\beta \xrightarrow{p} \frac{\sigma_Y}{2}(e_{Y,1}+e_{Y,2})$.
This estimate is \emph{contaminated} by the cross-cohort divergence:
its projection onto the true cohort-$k$ direction is
\begin{equation}
  \langle\hat\beta, e_{Y,k}\rangle \;=\; \frac{\sigma_Y}{2}(1+\rho_{YY}).
  \label{eq:betahat-Yk}
\end{equation}

\paragraph{Per-cohort batch estimate.}
Substituting $\hat\beta$ into the per-cohort residual average gives
\[
  \hat\gamma_k \;=\; \bar X_{\pi_k} - \alpha - \hat\beta\cdot\bar Y_{\pi_k}
  \;=\; \sigma_B\,(2k{-}3)\,e_B
      + \underbrace{\sigma_Y(\E[Y\mid B=k]-\tfrac12)\bigl(e_{Y,k}-\tfrac12(e_{Y,1}+e_{Y,2})\bigr)}_{\text{bias from }\rho_{YY}\ne 1}
      + O_p(n^{-1/2}).
\]
The bracketed vector simplifies to $\tfrac12(e_{Y,k}-e_{Y,3-k})$,
which is parallel to $e_{Y,1}-e_{Y,2}$ with norm
$\sqrt{2(1-\rho_{YY})}/2$.

\paragraph{Effect on post-correction label projection.}
Corrected data: $T_B(X_i)=X_i-\hat\gamma_{B_i}$. For a sample in
cohort $k$ with label $Y=1$,
\begin{align*}
  \langle T_B(X_i), e_{Y,k}\rangle
  &= \sigma_Y
    - \sigma_Y(\E[Y\mid B=k]-\tfrac12)\cdot\tfrac12
    \bigl(\langle e_{Y,k}, e_{Y,k}\rangle - \langle e_{Y,k},e_{Y,3-k}\rangle\bigr) \\
  &= \sigma_Y \bigl[1 - \tfrac14(2q-1)(1-\rho_{YY})\bigr]
      \qquad\text{for }k=1.
\end{align*}
The corresponding expression for $Y=0$ carries opposite sign, and
averaging across cohorts (with batch proportions 1/2 each) yields
the \emph{global} post-correction label discriminant
\begin{equation}
  \mathrm{disc}(\rho, \rho_{YY}, q)
  \;=\; \sigma_Y\Bigl[
     \tfrac12(1+\rho_{YY})
     \;-\; \rho^{2}
     \;-\; \tfrac12(2q-1)^{2}(1-\rho_{YY})
  \Bigr] \;+\; O_p(n^{-1/2}).
  \label{eq:disc-full}
\end{equation}
Equation~\eqref{eq:disc-full} is the precise identity underlying
Lemma~\ref{lem:retention}: it recovers the $(1-2\rho^2)$ form of
Theorem~1 in the special case $\rho_{YY}=1$, $q=1$ and reveals two
additional mechanisms (platform divergence $1-\rho_{YY}$ and
cohort-imbalance $2q-1$) that extend the flip condition to
\begin{equation}
  \mathrm{disc}<0
  \;\Leftrightarrow\;
  \rho^{2}
  \;+\;\tfrac12(2q-1)^{2}(1-\rho_{YY})
  \;>\;\tfrac12(1+\rho_{YY}).
  \label{eq:flip-extended}
\end{equation}
\paragraph{Extension to $K\ge 3$ cohorts.}
Let $\tilde B\in\R^{n\times(K-1)}$ be the centred batch-indicator
matrix; its column space
$\mathcal V_B=\mathrm{span}(\mu_{B=k}-\bar\mu)$ has generic rank
$K-1$. Let $U_B\in\R^{p\times(K-1)}$ have orthonormal columns
spanning $\mathcal V_B$. Define the \emph{alignment matrix}
\begin{equation}
  R \;=\; U_B^{\!\top}\,\mathrm{Col}\,[e_{Y,1}\;|\;\cdots\;|\;e_{Y,K}]
       \;\in\;\R^{(K-1)\times K},
  \label{eq:align-mat}
\end{equation}
i.e.\ the row of cosines between each column of $U_B$ and each
per-cohort label direction $e_{Y,k}$. Define the \emph{effective
alignment scalar}
\begin{equation}
  \rho^{2}_{\!\mathrm{eff}} \;=\; \frac{1}{K}\,\mathrm{tr}(R R^{\!\top}),
  \label{eq:rho-eff}
\end{equation}
and the generalised platform-divergence cross-correlation
\begin{equation}
  \rho_{YY}^{\!\mathrm{eff}} \;=\;
  \frac{2}{K(K-1)}\!\sum_{k<k'}\langle e_{Y,k}, e_{Y,k'}\rangle.
  \label{eq:rhoYY-eff}
\end{equation}

\begin{lemma}[General-$K$ retention identity]
\label{lem:retention-K}
Under the additive generative model
\eqref{eq:model-plat} extended to $K$ cohorts, the
global post-correction discriminant along the population label
direction $\bar e_Y=K^{-1}\sum_k e_{Y,k}$ satisfies
\begin{equation}
  \mathrm{disc}^{\!K}(\rho_{\mathrm{eff}},\rho_{YY}^{\!\mathrm{eff}},\{q_k\})
  \;=\; \sigma_Y\Bigl[
     \tfrac{1}{K}\!\sum_{k}\!\bigl(1+(K-1)\rho_{YY}^{\!\mathrm{eff}}\bigr)/K
     - \rho^{2}_{\mathrm{eff}}
     - c_{\!\boldsymbol q}(1-\rho_{YY}^{\!\mathrm{eff}})
  \Bigr] + O_p(n^{-1/2}),
  \label{eq:disc-K}
\end{equation}
where $c_{\!\boldsymbol q}=\tfrac{1}{K}\sum_k (2q_k-1)^2 /\,K$ is
the imbalance term.  The flip condition $\mathrm{disc}^K<0$ collapses
to Eq.~\eqref{eq:flip-extended} when $K=2$.
\end{lemma}

\begin{proof}[Proof (sketch, by induction on $K$).]
\emph{Base $K=2$.}  Direct calculation above (the three-parameter
identity).

\emph{Inductive step $K\to K+1$.}  Adjoin a new cohort $K+1$ with
centre $\mu_{B=K+1}$ and per-cohort direction $e_{Y,K+1}$.  The new
orthonormal column of $U_B$ is $u_{K+1}=(I-U_BU_B^{\!\top})(\mu_{K+1}
-\bar\mu)/\|\cdot\|$.  The updated $R$-matrix appends a new row
$u_{K+1}^{\!\top}[e_{Y,1}\;\cdots\;e_{Y,K+1}]$, increasing
$\mathrm{tr}(RR^{\!\top})$ by a single quadratic form in the new
cosines.  The weighted averages in
\eqref{eq:rho-eff}--\eqref{eq:rhoYY-eff} update by the usual
$K/(K+1)$ factor; substituting and collecting terms reproduces
\eqref{eq:disc-K} with $K\!+\!1$ in place of $K$.  The
$O_p(n^{-1/2})$ residual is uniform in $K$ because each new cohort
contributes an independent sub-Gaussian batch of size
$\Omega(n/(K+1))$ (Assumption~\ref{asm:n}).
\end{proof}

\begin{corollary}[$K$-generalised flip condition]
\label{cor:flipK}
$\mathrm{DIAL}>0$ iff
$\rho^{2}_{\!\mathrm{eff}} + c_{\!\boldsymbol q}(1-\rho_{YY}^{\!\mathrm{eff}})
> \tfrac{1}{K}\bigl(1+(K-1)\rho_{YY}^{\!\mathrm{eff}}\bigr)$.
For $K=2$ and $c_{\boldsymbol q}=(2q-1)^2/2$, this recovers
$\rho^2+\tfrac12(2q-1)^2(1-\rho_{YY}) > \tfrac12(1+\rho_{YY})$
as in Eq.~\eqref{eq:flip-extended}.
\end{corollary}

Corollary~\ref{cor:flipK} also implies: when platforms fully agree
($\rho_{YY}^{\!\mathrm{eff}}=1$) and cohort proportions are
balanced ($c_{\boldsymbol q}=0$), the flip requires
$\rho^{2}_{\!\mathrm{eff}}>1$, which is vacuous.  In the limit of
many cohorts $K\to\infty$ with balanced imbalance
($q_k=1/2+o(1)$), the RHS tends to $\rho_{YY}^{\!\mathrm{eff}}$, so
the flip requires $\rho^{2}_{\!\mathrm{eff}}>\rho_{YY}^{\!\mathrm{eff}}$
-- a clean geometric statement in the large-$K$ limit.

\section{Step (c): Flip at $\rho^{2}=\tfrac12$}

From Lemma~\ref{lem:retention}, the post-correction \emph{population}
discriminant along $w_Y$ is
\begin{equation}
  \mathrm{disc}(\rho) \;\propto\; 1-2\rho^{2}.
  \label{eq:disc}
\end{equation}
Thus:
\begin{itemize}
\item If $\rho^2<\tfrac12$, $\mathrm{disc}(\rho)>0$ and the population
  classifier is correctly aligned with $w_Y$.
\item If $\rho^2>\tfrac12$, $\mathrm{disc}(\rho)<0$ and the population
  classifier is anti-aligned with $w_Y$.
\item Equality $\rho^2=\tfrac12$ is a measure-zero boundary.
\end{itemize}
\TODO{The constant $\tfrac12$ is exact for two balanced cohorts
($K=2$, $|\pi_1|=|\pi_2|$). For $K\ge 3$ or imbalanced cohorts, the
boundary is $c_\ast(K,\boldsymbol{\pi})$; equations~\eqref{eq:disc} and
Lemma~\ref{lem:retention} should be recomputed with the exact
Gram--Schmidt coefficients to pin down $c_\ast$.}

\section{Step (d): Finite-sample consistency of $w_{\text{trained}}$}

Under Assumption~\ref{asm:n} and the sub-Gaussian residual hypothesis,
standard ridge-regression / M-estimator concentration (Hsu, Kakade,
Zhang 2014; Negahban et al.\ 2012) gives
\begin{equation}
  w_{\text{trained}} \;=\; \eta \cdot \mathrm{disc}(\rho)\, w_Y
    + O_{\!p}\!\bigl(\sqrt{p/n}\bigr)
\label{eq:wtr}
\end{equation}
for some $\eta>0$ depending on the loss curvature.
\TODO{Fill in the exact concentration inequality; the hinge-loss case
requires Talagrand's contraction lemma and a VC-entropy bound on the
feature space.}

\paragraph{Sign argument.}
From~\eqref{eq:wtr}, $\mathrm{sign}\langle w_{\text{trained}}, w_Y\rangle
= \mathrm{sign}(1-2\rho^{2})$ with probability $1-o(1)$ as $n\to\infty$.

\section{Step (e): Translation to AUC}

For a held-out cohort, conditional on $w_{\text{trained}} = \eta'\, w_Y$
(in the $\rho^2<\tfrac12$ regime) or $w_{\text{trained}}=-\eta'\,w_Y$
(in the $\rho^2>\tfrac12$ regime), the Mann--Whitney statistic of the
linear score $w_{\text{trained}}^{\!\top}X^{\text{test}}$ versus the
label satisfies
\begin{equation}
   \lim_{n\to\infty}\mathrm{AUC}_{\text{post}} \;=\;
   \begin{cases}
     \Phi\!\bigl(\tfrac12\lVert\mu_1-\mu_0\rVert/\sigma\bigr) & \rho^2<\tfrac12,\\[2pt]
     1-\Phi\!\bigl(\tfrac12\lVert\mu_1-\mu_0\rVert/\sigma\bigr) & \rho^2>\tfrac12,
   \end{cases}
\end{equation}
where $\Phi$ is the standard normal CDF and $\sigma^2$ is the residual
variance along $w_Y$.  In particular, the two regimes have AUC values
symmetric around $1/2$.  The \emph{direction-invariant} measurement
$\max(\mathrm{AUC},1-\mathrm{AUC})$ is identical in both regimes while
the raw AUC flips, proving both parts of Theorem~\ref{thm:flip} and
Corollary~\ref{cor:dial}.

\section{Gap summary for user verification}

\begin{enumerate}
\item \TODO{Step (b) algebraic identity: derive~\eqref{eq:retention}
  rigorously for $K=2$ balanced cohorts, then generalise to $K\ge 3$.}
\item \TODO{Step (c) constant: confirm that the threshold
  $\rho^{\star 2}=\tfrac12$ is sharp for $K=2$; explicitly compute
  $c_\ast(K,\boldsymbol{\pi})$ for imbalanced cohorts.}
\item \TODO{Step (d) concentration: give a finite-sample bound with
  explicit dependence on $p/n$ and on the loss curvature.}
\item \TODO{Check the universal constant $\tfrac12$ does not become
  $\tfrac12(1+\mathrm{Var}(Y))$ after accounting for the label covariate
  contribution; Johnson et al.\ 2007 eq. 3 hints this possibility.}
\item \TODO{Relax the isotropic-residual assumption to general
  $\Sigma=U\Lambda U^{\!\top}$; the flip condition becomes
  $\rho^{\!\top}\Lambda^{-1}\rho>\tfrac12 \rho^{\!\top}\Lambda^{-1}\rho$
  which simplifies only for specific $\Sigma$.}
\end{enumerate}

\section{Corollary 1 (proof)}

\begin{proof}[Proof of Corollary~\ref{cor:dial}]
By Theorem~\ref{thm:flip},
$\mathrm{AUC}_{\text{post}}\to\Phi(\cdot)$ or $1-\Phi(\cdot)$ deterministically
in the two regimes, and DIAL$=\max(\mathrm{AUC},1-\mathrm{AUC})-\mathrm{AUC}$
is identically $0$ in the non-flip regime and strictly positive in the
flip regime. Equivalence with $\rho^2>\tfrac12$ follows immediately.
\end{proof}

\end{document}
